\section{文獻回顧}\label{sec:lit}

本節的回顧依一條線索展開：小人口死亡率估計的困難出在哪裡、
死亡率文獻以哪兩類作法應對、這兩類作法各自所依的一般統計理論為何，
以及它們共同留下的缺口。前兩小節屬死亡率文獻，第三小節是本文病灶的直接
上游，第四小節為本文推導所借用的一般統計理論與本文結果所要接上的精算實務
標準，末一小節交代後續各節的安排。各條依其發展順序書寫，
並於末尾說明該脈絡所留下的缺口、其原始文獻自己記載的限制，
以及本文自該脈絡所借用的性質。

\subsection{Lee-Carter 模型、其變體，與小人口下的參數偏誤}\label{subsec:lc}

最早把年齡與時間分離的雙線性死亡率模型由 \citet{lee1992} 提出。
該文以奇異值分解配適 $\log m_{x,t}=\alpha_x+\beta_x\kappa_t$，
再以時間序列模型外推 $\kappa_t$；由於形式簡單、參數可解釋、
且外推只需處理一個單變數序列，很快成為各國死亡率預測的基準模型，
亦為聯合國與多國統計機構所採用。

在應用的過程中，三項技術問題陸續被指出，而每一項都各自引出一條補救。
首先是配適末年與觀測值的落差。\citet{leemiller2001} 以多國資料評估原始
作法的預測表現，發現第二階段重解 $\kappa_t$ 時若以總死亡數對齊，
配適末年的死亡率與觀測值仍有系統性差距，因而建議改為對齊當年的零歲平均
餘命，此即後續文獻所稱的 Lee--Miller 變體。
其次是誤差變異的設定。\citet{brouhns2002} 指出奇異值分解隱含各年齡的誤差
變異數相等，而死亡數為計數資料、其變異隨水準而變，兩者不符；
該文因而改以 Poisson 對數雙線性模型直接配適計數，
把估計由對數尺度的最小平方換為計數尺度的最大概似。
再者是 $\hat\beta_x$ 的不平滑。\citet{currie2004} 發現逐年齡估計的
$\hat\beta_x$ 在相鄰年齡之間劇烈起伏，而真實的年齡型態應為平滑，
該文因而把 \citet{eilers1996} 的懲罰 B-樣條引入死亡率配適；
\citet{delwarde2007} 則在 Poisson 對數概似上直接加差分懲罰，
使修勻與配適在同一個目標函數內完成。

上述補救本身又引出新的問題。\citet{djeundje2010} 發現以年齡與年份分類的
死亡資料常有過度離散，其變異數大於均值，故 Poisson 假設雖較常態為佳
卻仍不足，該文因而放寬為「變異數與均值成比例」的擬似概似設定。
\citet{currie2013} 進一步把這一類帶約束與懲罰的廣義線性模型寫成一般的
計算框架，使各種變體可在同一套程式內實作。
模型本身的擴充亦同時進行，\citet{renshaw2003} 加入年齡別的增強項，
\citet{camarda2021} 則把修勻、分解與預測整合於同一套平滑架構。
\citet{rabbi2021} 另指出修勻的選擇會改動預測所隱含的壽命離散度，
故修勻並非只影響參數的外觀。

這些變體隨後被系統地評比。\citet{booth2006} 以多國資料比較各變體的預測
表現，\citet{shang2011} 進一步把十種主成分型作法並列，
同時比較點預測與區間預測，指出各變體之間的差異在區間上比在點上更明顯。
\citet{hunt2020} 則系統刻畫了年齡--時期類模型的識別性問題，
指出若不施加額外約束，參數的個別數值並無唯一解。
值得一併指出的是，Lee-Carter 第一階段的奇異值分解在函數型資料的語彙中即
為主成分分析：\citet{hyndman2007} 以穩健的函數型主成分分析處理戰爭與疫情
等離群年份，\citet{hormann2015} 則把主成分推廣到容許時間相依的動態版本。
\citet{basellini2023} 在模型提出三十年後做了整體回顧，
並把小人口列為尚未解決的課題之一。

小人口下 Lee-Carter 估計的核心現象是年齡參數的低估。既有文獻的共同發現是，
人數減少時 $\hat\alpha_x$ 與 $\hat\beta_x$ 會朝同一方向偏移，
而該偏移在重複抽樣中方向一致，故屬偏誤而非變異。
\citet{chen2017} 最早把人口規模獨立為研究的對象，以英格蘭與威爾斯的資料
指出隨機死亡率模型的參數不確定性與人口規模直接相關，在小人口下不可忽略，
並區分了來自母體規模的抽樣變異與來自模型的參數變異。
\citet{wang2018} 以模擬給出該偏移的量級，指出 Lee-Carter 的年齡參數
$\hat\alpha_x$ 與 $\hat\beta_x$ 在人數減少時同向偏移，
且在人數二十萬以下時特別明顯。
\citet[Figure 1]{yue2019} 把兩個參數分開，發現 $\hat\alpha_x$ 的偏誤恆為正
且量級遠大於 $\hat\beta_x$，而後者在人數超過十萬時平均而言已接近零。

國內的研究亦指向同一現象，且其規模設定更貼近鄉鎮市區。
陳政勳與余清祥（2010）探討臺北市、雲嘉兩縣與澎湖縣的小區域人口推估，
王信忠等（2012）以區塊拔靴法發現南港區（單一性別約五萬人）的平均餘命
預測區間寬度超過十歲，王信忠等（2017）編製規模與鄉鎮市區相當的臺灣原住民
生命表，余清祥等（2021）以年輪變動比處理小區域人口推估，
余清祥等（2025）則改以標準化死亡率與平均餘命的線性關係繞過生命表的編算。

這條脈絡至此已完成現象的提出、量化與分解，而它所留下的缺口在於\textbf{描述
的形式}：上述各文對偏誤的陳述一律是模擬所得的現象，
未見有研究上溯至該偏誤的生成環節，亦未見其量的解析形式。
缺乏解析形式的直接後果是，使用者無從在配適之前判斷某一組資料會偏多少，
只能逐案以模擬估計；而模擬須先指定真值，這在實際的小區域應用中並不可得。
以下兩小節回顧文獻對此的兩類應對：一類設法補進外部資訊，
另一類設法改動估計的機制。
